% Macro safety net for standalone extraction.
\providecommand{\BRST}{\mathrm{BRST}}
\providecommand{\BRSTGauged}{\mathrm{BRSTGaugedChirAlg}}
\providecommand{\ClaimStatusComputed}{\textnormal{[computed]}}

% Regime III: universal, freely generated W_3 algebra over C(c),
% followed by specialization away from the Shapovalov divisor.
\section{The Zamolodchikov \texorpdfstring{$\mathcal W_3$}{W3}
normalization}
\label{sec:w3-composite-complete}

The universal $\mathcal W_3$ vertex algebra is generated by a
Virasoro field $T$ of weight~$2$ and a Virasoro-primary field $W$ of
weight~$3$.  We use
\[
T(z)=\sum_{n\in\mathbb Z}L_nz^{-n-2},\qquad
W(z)=\sum_{n\in\mathbb Z}W_nz^{-n-3},
\]
and fix the scale of $W$ by
\[
W(z)W(w)\sim \frac{c/3}{(z-w)^6}
              +\frac{2T(w)}{(z-w)^4}+\cdots.
\]
This convention determines every coefficient below.  In particular,
the coefficient of the weight-$4$ quasi-primary in the second-order
pole is $32/(22+5c)$.  The coefficient $16/(22+5c)$ belongs to the
mode commutator.  The contour calculation in
Proposition~\ref{prop:w3-ope-mode-normalization} explains the factor
of two.

\subsection{The level-four Virasoro quasi-primary}

\begin{definition}[The composite field $\Lambda$]
\label{def:lambda-complete}
\index{Lambda@$\Lambda$ (composite field)|textbf}
\index{W3 algebra@$\mathcal W_3$!composite field Lambda@composite field $\Lambda$|textbf}
The level-$4$ Virasoro quasi-primary is
\begin{equation}\label{eq:lambda-qp}
\Lambda={:}TT{:}-\frac{3}{10}\partial^2T.
\end{equation}
Its state and mode expansions are
\begin{equation}\label{eq:lambda-state}
|\Lambda\rangle
=L_{-2}^2|0\rangle-\frac35L_{-4}|0\rangle,
\end{equation}
\begin{equation}\label{eq:lambda-mode-correct}
\Lambda_n
=\sum_{p\in\mathbb Z}{:}L_pL_{n-p}{:}
-\frac{3}{10}(n+2)(n+3)L_n,
\end{equation}
where Virasoro vacuum normal ordering means
\begin{equation}\label{eq:vir-mode-normal-ordering}
{:}L_pL_q{:}
=
\begin{cases}
L_pL_q,&p\leq-2,\\
L_qL_p,&p\geq-1.
\end{cases}
\end{equation}
\end{definition}

\begin{theorem}[Level-four calculation and normalized $WW$ coupling;
\ClaimStatusProvedElsewhere]
\label{thm:lambda-coefficients-derivation}
In the Zamolodchikov normalization, the following identities hold:
\begin{align}
L_1|\Lambda\rangle&=0,\label{eq:lambda-L1}\\
L_2|\Lambda\rangle&=\frac{5c+22}{5}L_{-2}|0\rangle,
\label{eq:lambda-L2}\\
\langle\Lambda|\Lambda\rangle&=\frac{c(5c+22)}{10},
\label{eq:lambda-norm}\\
W_{(1)}W\big|_{\mathbb C\Lambda}
&=\frac{32}{22+5c}\,\Lambda.
\label{eq:w3-lambda-ope-coupling}
\end{align}
Equivalently,
\begin{equation}\label{eq:T-lambda-ope}
T(z)\Lambda(w)\sim
\frac{(5c+22)T(w)/5}{(z-w)^4}
+\frac{4\Lambda(w)}{(z-w)^2}
+\frac{\partial\Lambda(w)}{z-w}.
\end{equation}
The final identity is the Zamolodchikov Jacobi solution
\cite{Zamolodchikov} in the conventions reviewed by
Bouwknegt--Schoutens~\cite[\S3.1]{Bouwknegt-Schoutens}; their
parameter is $b^2=16/(22+5c)$ and the second-order pole contains
$2b^2\Lambda$.  The first three identities
follow directly from the Virasoro relations.
\end{theorem}

\begin{proof}
Write
$|v_\beta\rangle=L_{-2}^2|0\rangle+2\beta L_{-4}|0\rangle$.
The Virasoro commutator gives
\[
L_1L_{-2}^2|0\rangle=3L_{-3}|0\rangle,
\qquad
L_1L_{-4}|0\rangle=5L_{-3}|0\rangle.
\]
Thus $L_1|v_\beta\rangle=(3+10\beta)L_{-3}|0\rangle$,
and quasi-primarity fixes $\beta=-3/10$.
Next,
\[
L_2\left(L_{-2}^2-\frac35L_{-4}\right)|0\rangle
=\left(c+8-\frac{18}{5}\right)L_{-2}|0\rangle
=\frac{5c+22}{5}L_{-2}|0\rangle.
\]
Finally, the level-$4$ Shapovalov products are
\[
\langle L_{-2}^2|L_{-2}^2\rangle=\frac{c(c+8)}2,
\quad
\langle L_{-4}|L_{-2}^2\rangle=3c,
\quad
\langle L_{-4}|L_{-4}\rangle=5c.
\]
Substitution yields
\[
\langle\Lambda|\Lambda\rangle
=\frac{c(c+8)}2-\frac65(3c)+\frac9{25}(5c)
=\frac{c(5c+22)}{10}.
\]
The Jacobi identity for $(T,W,W)$ and $(W,W,W)$ then fixes the
coefficient in~\eqref{eq:w3-lambda-ope-coupling}; this is the
standard $\mathcal W_3$ presentation of
Zamolodchikov~\cite{Zamolodchikov}.
\end{proof}

\subsection{The complete singular $WW$ product}

\begin{theorem}[Zamolodchikov $W$--$W$ OPE;
\ClaimStatusProvedElsewhere]
\label{thm:w-w-ope-complete}
Over $\mathbb C(c)$, the complete singular part of the $W$--$W$ OPE
is
\begin{align}
W(z)W(w)
&\sim \frac{c/3}{(z-w)^6}
+\frac{2T(w)}{(z-w)^4}
+\frac{\partial T(w)}{(z-w)^3}\nonumber\\
&\quad
+\frac{\frac{3}{10}\partial^2T(w)
       +\frac{32}{22+5c}\Lambda(w)}{(z-w)^2}
+\frac{\frac{1}{15}\partial^3T(w)
       +\frac{16}{22+5c}\partial\Lambda(w)}{z-w}.
\label{eq:w-w-ope-full}
\end{align}
Here $\Lambda$ is the field in
Definition~\ref{def:lambda-complete}.  The derivative coefficient in
the first-order pole is one half of the quasi-primary coefficient in
the second-order pole, as required by translation covariance.
\end{theorem}

\begin{proof}
The leading pole and the coefficient of $T$ fix the scale of $W$.
Virasoro covariance supplies the descendant coefficients
$1$, $3/10$, and $1/15$ in the $T$ conformal family.  The only new
level-$4$ quasi-primary is $\Lambda$, and the Jacobi identity fixes
its coefficient to $32/(22+5c)$
\cite{Zamolodchikov,Bouwknegt-Schoutens}.  Translation covariance
then supplies $16/(22+5c)$ on $\partial\Lambda$.
\end{proof}

\begin{proposition}[OPE-to-mode normalization;
\ClaimStatusProvedHere]
\label{prop:w3-ope-mode-normalization}
Suppose the $\Lambda$ family appears in a weight-three self-OPE as
\[
\frac{A\Lambda(w)}{(z-w)^2}
+\frac{A\,\partial\Lambda(w)/2}{z-w}.
\]
Its contribution to the mode commutator is
\begin{equation}\label{eq:ope-to-mode-half}
[W_m,W_n]\big|_\Lambda
=\frac{A}{2}(m-n)\Lambda_{m+n}.
\end{equation}
Consequently the OPE coefficient
$A=32/(22+5c)$ in~\eqref{eq:w-w-ope-full} gives the standard mode
coefficient $16(m-n)/(22+5c)$.
\end{proposition}

\begin{proof}
The double contour defining $[W_m,W_n]$ uses the factors
$z^{m+2}w^{n+2}$.  The second-order pole contributes
$A(m+2)\Lambda_{m+n}$.  Since
$\partial\Lambda(w)=\sum_p(-p-4)\Lambda_pw^{-p-5}$, the first-order
pole contributes $-A(m+n+4)\Lambda_{m+n}/2$.  Their sum is
$A(m-n)\Lambda_{m+n}/2$.
\end{proof}

\begin{corollary}[Zamolodchikov mode algebra;
\ClaimStatusProvedElsewhere]
\label{cor:w3-mode-commutator}
The $W$ modes satisfy
\begin{align}
[W_m,W_n]
&=\frac{c}{360}m(m^2-1)(m^2-4)\delta_{m+n,0}\nonumber\\
&\quad +(m-n)
\left[
\frac{(m+n+3)(m+n+2)}{15}
-\frac{(m+2)(n+2)}6
\right]L_{m+n}\nonumber\\
&\quad+\frac{16}{22+5c}(m-n)\Lambda_{m+n}.
\label{eq:w3-mode-commutator}
\end{align}
\end{corollary}

\subsection{Highest-weight check}

\begin{proposition}[Action of $\Lambda_0$;
\ClaimStatusProvedHere]
\label{prop:lambda-zero-highest-weight}
Let $|h,w\rangle$ be a $\mathcal W_3$ highest-weight vector.  Then
\begin{equation}\label{eq:lambda-zero-eigenvalue}
\Lambda_0|h,w\rangle
=\left(h^2+\frac h5\right)|h,w\rangle.
\end{equation}
\end{proposition}

\begin{proof}
Formula~\eqref{eq:lambda-mode-correct} and the normal ordering
in~\eqref{eq:vir-mode-normal-ordering} give
\[
\sum_p{:}L_pL_{-p}{:}|h,w\rangle
=(L_0^2+L_1L_{-1})|h,w\rangle
=(h^2+2h)|h,w\rangle.
\]
The derivative correction is $-(9/5)L_0$, which gives
\eqref{eq:lambda-zero-eigenvalue}.
\end{proof}

\begin{theorem}[Level-one Shapovalov matrix;
\ClaimStatusProvedHere]
\label{thm:w3-kac-level1}
In the ordered basis
$\{L_{-1}|h,w\rangle,W_{-1}|h,w\rangle\}$, the level-one
Shapovalov matrix is
\begin{equation}\label{eq:w3-gram-level1}
M_1(c,h,w)=
\begin{pmatrix}
2h&3w\\[2pt]
3w&-\frac h5+\frac{32}{22+5c}\left(h^2+\frac h5\right)
\end{pmatrix}.
\end{equation}
For $h\ne0$, a level-one singular vector exists precisely on
\begin{equation}\label{eq:w3-null-curve}
9w^2(22+5c)=2h^2(32h+2-c),
\end{equation}
and may be chosen as
\begin{equation}\label{eq:w3-null-vector}
\left(W_{-1}-\frac{3w}{2h}L_{-1}\right)|h,w\rangle.
\end{equation}
\end{theorem}

\begin{proof}
The Virasoro-primary relation gives
$[L_1,W_{-1}]=3W_0$, so the first row is $(2h,3w)$.
Equations~\eqref{eq:w3-mode-commutator} and
\eqref{eq:lambda-zero-eigenvalue} give
\[
[W_1,W_{-1}]|h,w\rangle
=\left[-\frac h5+
\frac{32}{22+5c}\left(h^2+\frac h5\right)\right]|h,w\rangle.
\]
This proves~\eqref{eq:w3-gram-level1}.  Its determinant vanishes
exactly when~\eqref{eq:w3-null-curve} holds.  The first-row kernel
then gives~\eqref{eq:w3-null-vector}.
\end{proof}

\subsection{Scalar exchange ratios and their scope}

Set
\begin{equation}\label{eq:w3-two-couplings}
D:=5c+22,
\qquad
A_{\mathrm{OPE}}:=\frac{32}{D},
\qquad
\beta_{\mathrm{mode}}:=\frac{16}{D}.
\end{equation}
The distinction between $A_{\mathrm{OPE}}$ and
$\beta_{\mathrm{mode}}$ determines the factor-of-four distinction in
a squared exchange coefficient.

\begin{proposition}[Rank-one $\Lambda$-exchange ratios;
\ClaimStatusComputed]
\label{prop:w3-lambda-exchange-ratios}
Away from $cD=0$, inversion of the one-dimensional Gram line
$\mathbb C\Lambda$ gives
\begin{align}
\frac{1}{\langle\Lambda|\Lambda\rangle}
&=\frac{10}{cD},\label{eq:lambda-inverse-norm}\\
\frac{A_{\mathrm{OPE}}}{\langle\Lambda|\Lambda\rangle}
&=\frac{320}{cD^2},&
\frac{A_{\mathrm{OPE}}^2}{\langle\Lambda|\Lambda\rangle}
&=\frac{10240}{cD^3},\label{eq:lambda-ope-ratios}\\
\frac{\beta_{\mathrm{mode}}}{\langle\Lambda|\Lambda\rangle}
&=\frac{160}{cD^2},&
\frac{\beta_{\mathrm{mode}}^2}{\langle\Lambda|\Lambda\rangle}
&=\frac{2560}{cD^3}.
\label{eq:lambda-mode-ratios}
\end{align}
The second line is OPE-normalized; the third is mode-normalized.
Each line is a rank-one exchange ratio.  A full quartic shadow
requires a specified extraction map, its symmetry factors, and the
complete intermediate-state Gram block.
\end{proposition}

\begin{proof}
Insert~\eqref{eq:lambda-norm} and
\eqref{eq:w3-two-couplings}.  The final sentence records the domain
of this one-dimensional calculation.
\end{proof}

\begin{remark}[Weight-six interface]
\label{rem:w3-weight-six-interface}
For the universal freely generated vacuum module, the character
\[
\prod_{n\ge2}(1-q^n)^{-1}\prod_{n\ge3}(1-q^n)^{-1}
\]
has coefficients $\dim V_5=4$ and $\dim V_6=8$.  On the generic
translation-injective locus, the quotient
$V_6/L_{-1}V_5$ therefore has dimension~$4$.  The regular product
$W_{(-1)}W={:}WW{:}$ determines one vector in this quotient.  A
weight-six Gram determinant requires a chosen basis of the full
four-dimensional quotient and a direct Shapovalov computation.
Accordingly, the factors of such a determinant and every quartic
formula using them remain an explicit finite calculation rather than
an inference from the singular OPE.
\end{remark}

\begin{remark}[Bar and modular interface]
\label{rem:regular-term-bar}
The OPE theorem fixes the local vertex-algebra structure.  A chiral
bar differential additionally requires the chosen curve, completion,
factorization product, and twisting morphism.  A modular scalar
requires the genus-one trace or curvature construction.  Thus
$c/2$, $c/3$, and the ratios in
Proposition~\ref{prop:w3-lambda-exchange-ratios} are local
normalization data; their promotion to bar curvature, Koszul duality,
or a higher-genus shadow is governed by the corresponding named
hypothesis package in Chapter~\ref{chap:w-algebra-koszul}.
\end{remark}
